% Standalone Beamer deck — open in the file tree and click Compile.
\documentclass[11pt,aspectratio=169]{beamer}

\usetheme{Madrid}
\usecolortheme{default}
\setbeamertemplate{navigation symbols}{}
\setbeamertemplate{footline}[frame number]

\usepackage[T1]{fontenc}
\usepackage[utf8]{inputenc}
\usepackage{amsmath,amssymb}
\usepackage{booktabs}
\usepackage{hyperref}

\title[Multi-file \LaTeX]{Reproducible Research in Practice}
\subtitle{Paper, supplement \& slides from one project}
\author{First Author \and Second Author}
\institute{University of Example}
\date{\today}

\begin{document}

\begin{frame}
  \titlepage
\end{frame}

\begin{frame}{Why multiple PDFs in one project?}
  \begin{itemize}
    \item \textbf{Main paper} — journal or conference article (\texttt{main.tex})
    \item \textbf{Supplementary PDF} — extended tables, proofs (\texttt{supplement.tex})
    \item \textbf{Slides} — talk deck for the same work (\texttt{slides.tex})
    \item \textbf{Shared assets} — one \texttt{references.bib}, figures, data files
  \end{itemize}
  \vfill
  \small One \textbf{Compile} button: open the file you want, then compile.
\end{frame}

\begin{frame}{Workflow}
  \begin{enumerate}
    \item Edit \texttt{main.tex}, \texttt{supplement.tex}, or \texttt{slides.tex}
    \item Banner shows \emph{Compile builds this file} for standalone \texttt{.tex}
    \item Server runs pdfLaTeX + BibTeX when needed
    \item Preview updates for \textbf{that} document only
  \end{enumerate}
\end{frame}

\begin{frame}{Results at a glance}
  \begin{table}
    \centering
    \begin{tabular}{lcc}
      \toprule
      Method & Accuracy & Time (s) \\
      \midrule
      Baseline & 0.82 & 118 \\
      Proposed & \textbf{0.91} & \textbf{43} \\
      \bottomrule
    \end{tabular}
  \end{table}
  \vfill
  \footnotesize Details in the supplement; see~\cite{example2024methods}.
\end{frame}

\begin{frame}{References}
  \footnotesize
  \bibliographystyle{IEEEtran}
  \bibliography{references}
\end{frame}

\begin{frame}
  \centering
  \Large Questions?
  \vfill
  \small \url{https://example.org}
\end{frame}

\end{document}
